% DRAFT text for sn-article_v1.tex (Abstract, Contributions, Table 2, Discussion, Conclusion).
% Fragments, not compilable alone. Every number is from the 5-seed re-run and the HGT investigation:
%   docs/EXPERIMENT_LOG_semantic_edge_investigation.md  ("Baseline fairness re-run (#16)" and "HGT investigation"),
%   results/*_5seed.csv, results/hgt_controls_5seed.csv, docs/LATENCY.md.
% The last block ("AUTHOR NOTES") lists statements elsewhere in the paper that this text contradicts. NOT FOR SUBMISSION.

%%%%%%%%%%%%%%%%%%%%%%%%%%%%%%%%%%%%%%%%%%%%%%%%%%%%%%%%%%%%%%%%%%%%%%%%%%%%%%%%%%%%%%%%%%%
%% ABSTRACT (replaces \abstract{...})
%%%%%%%%%%%%%%%%%%%%%%%%%%%%%%%%%%%%%%%%%%%%%%%%%%%%%%%%%%%%%%%%%%%%%%%%%%%%%%%%%%%%%%%%%%%
\abstract{Signature-based web application firewalls and sequence models treat an HTTP request as a flat token stream and are evaded by obfuscation that fragments keywords. We represent each request as a directed multi-relational token graph with sequential, skip and rule-derived semantic edges ($E_{sem}$), the last of which reconstructs keyword pairs even when they are split by comments or whitespace, and classify it with a GATv2 network (LOGWAT). We re-train seven baselines under a protocol identical to LOGWAT's and find that the standard test split cannot discriminate architectures (all $\geq$ 99.9\% accuracy) and that the margins reported in earlier comparisons do not reproduce; we therefore compare models on a held-out obfuscation matrix and an external dataset. There, LOGWAT is statistically indistinguishable from GCN, GraphSAGE, GIN and a token-level TextCNN on the held-out matrix (78.0 $\pm$ 4.5 of 90 correct versus 76--83), fails its comment-splitting cell at every seed, and attains the best mean weighted F1 on the external dataset (0.796 $\pm$ 0.030). We trace the failure to feature dilution: under one shared aggregation the four semantic edges of a split payload are outnumbered about 35:1 by n-gram edges. A relation-typed aggregator (HGT) solves the cell at every seed; controlled retraining of the same architecture without relation information (solved in 2 of 10 runs versus 5 of 5) and edits to the four semantic edges attribute this to edge typing rather than to attention form or capacity. The benefit is bounded by the semantic-edge builder: it transfers to new splitting styles only when the builder emits a semantic edge (25/30 versus 7/30 for GATv2; 0/15 when it does not), and typed HGT is seed-unstable on external data (weighted F1 0.33--0.82), which we trace to very short requests absent from training. LOGWAT itself runs at about 2~ms per request.}

%%%%%%%%%%%%%%%%%%%%%%%%%%%%%%%%%%%%%%%%%%%%%%%%%%%%%%%%%%%%%%%%%%%%%%%%%%%%%%%%%%%%%%%%%%%
%% CONTRIBUTIONS (replaces the itemize at the end of the Introduction)
%%%%%%%%%%%%%%%%%%%%%%%%%%%%%%%%%%%%%%%%%%%%%%%%%%%%%%%%%%%%%%%%%%%%%%%%%%%%%%%%%%%%%%%%%%%
The main contributions of this work are summarized as follows:
\begin{itemize}
    \item We propose a multi-relational behavioral attack graph (BAG) for HTTP requests with sequential, skip and semantic edges, where $E_{sem}$ is built by a keyword-pair matcher that reconstructs keywords split by comment or whitespace noise, so that different obfuscations of one payload map to the same semantic edge.
    \item We provide a like-for-like comparison of LOGWAT (GATv2) with seven baselines re-trained under the same protocol, split, seeds and evaluation. It shows that a standard test split saturates for every model and cannot rank architectures, and it identifies the held-out obfuscation matrix and an external dataset as the settings in which they differ.
    \item We diagnose LOGWAT's one systematic held-out failure (comment splitting) as dilution of the rare semantic edges by dense n-gram edges under a shared aggregation, and isolate its remedy: a relation-typed aggregation path, established by controlled retraining and inference-time edge edits rather than by comparing architectures alone.
    \item We characterize the limits of that remedy: it is bounded by the semantic-edge builder, produces false alarms on benign text containing keyword pairs, and is seed-unstable on very short requests that the training data do not cover.
    % [BIA bullet: keep only if supported by docs/BIA_DECISION_LOG.md; not evaluated in this investigation.]
\end{itemize}

%%%%%%%%%%%%%%%%%%%%%%%%%%%%%%%%%%%%%%%%%%%%%%%%%%%%%%%%%%%%%%%%%%%%%%%%%%%%%%%%%%%%%%%%%%%
%% TABLE 2 (replaces Table~\ref{tab:method_comparison_extended}); mean $\pm$ std over seeds 42--46
%%%%%%%%%%%%%%%%%%%%%%%%%%%%%%%%%%%%%%%%%%%%%%%%%%%%%%%%%%%%%%%%%%%%%%%%%%%%%%%%%%%%%%%%%%%
\begin{table}[t]
\centering
\caption{Baselines re-trained with LOGWAT's exact protocol (same graphs or tokens, split, loss, optimizer, early stopping, default learning rate; five seeds). Test = frozen test split accuracy (\%); Held-out = correct of 90 rows over 9 obfuscation cells; External = weighted F1 on the external dataset (30{,}677 requests); Latency = end-to-end mean, batch size 1, CPU (ms), measured in one run. Latency of the sequence models was not measured.}
\label{tab:method_comparison_extended}
\footnotesize
\setlength{\tabcolsep}{4pt}
\begin{tabular}{|l|c|c|c|c|}
\hline
\textbf{Method} & \textbf{Test} & \textbf{Held-out /90} & \textbf{External wF1} & \textbf{Latency} \\
\hline
TextCNN    & 100.000 & 82.6 $\pm$ 5.7  & 0.509 $\pm$ 0.046 & -- \\
Bi-LSTM    & 100.000 & 71.4 $\pm$ 17.8 & 0.444 $\pm$ 0.079 & -- \\
Stack-LSTM & 99.938  & 50.6 $\pm$ 16.0 & 0.392 $\pm$ 0.035 & -- \\
GCN        & 99.991  & 80.0 $\pm$ 7.1  & 0.745 $\pm$ 0.057 & 1.41 \\
GraphSAGE  & 99.991  & 78.0 $\pm$ 4.5  & 0.715 $\pm$ 0.064 & 1.14 \\
GIN        & 99.991  & 76.0 $\pm$ 5.5  & 0.692 $\pm$ 0.075 & 1.15 \\
HGT (typed edges) & 99.991 & \textbf{90.0 $\pm$ 0.0} & 0.535 $\pm$ 0.218 & 5.60 \\
\textbf{LOGWAT (GATv2)} & 99.981 & 78.0 $\pm$ 4.5 & \textbf{0.796 $\pm$ 0.030} & 1.79 \\
\hline
\end{tabular}
\end{table}
% RoBERTa / CodeBERT (single training run, seed 42; not in the table): held-out 60/90 and 70/90; external wF1 0.778 and 0.533.

%%%%%%%%%%%%%%%%%%%%%%%%%%%%%%%%%%%%%%%%%%%%%%%%%%%%%%%%%%%%%%%%%%%%%%%%%%%%%%%%%%%%%%%%%%%
%% DISCUSSION (replaces \subsection{Discussion} and its three bullets)
%%%%%%%%%%%%%%%%%%%%%%%%%%%%%%%%%%%%%%%%%%%%%%%%%%%%%%%%%%%%%%%%%%%%%%%%%%%%%%%%%%%%%%%%%%%
\subsection{Discussion}

\paragraph{The test split does not rank architectures.}
When every baseline is re-trained on the current dataset with LOGWAT's own protocol, all eight models reach at least 99.9\% test accuracy (Table~\ref{tab:method_comparison_extended}); each baseline lies 4--10 points above its earlier figure, and its margin to LOGWAT falls within $\pm$0.04 points (about two of 4{,}215 samples). The 2--8-point advantage reported in our earlier comparison therefore does not reproduce, and the old table should be read as superseded rather than corrected. We do not claim that the earlier baselines were mis-configured: their configurations cannot be recovered and the dataset changed as well. We attribute the saturation to the split rather than to equal capability: 25.2\% of test requests are byte-identical to a training request, and the SQLi class is drawn almost entirely from a pool of 41 payload templates. Checkpoints for every graph and sequence model, LOGWAT included, are selected on test-split accuracy, so these numbers are optimistic in the same way for all.

\paragraph{Held-out and external results.}
Discrimination comes from two settings that the training procedure never optimizes. On the held-out matrix, LOGWAT (78.0 $\pm$ 4.5 of 90) cannot be told apart from GCN, GraphSAGE, GIN or TextCNN (76--83), and its loss is concentrated in one cell, comment splitting, which it fails at every seed (0/50). On the external dataset, LOGWAT has the best mean weighted F1 (0.796 $\pm$ 0.030), clearly above the from-scratch sequence models ($\leq$0.51) and modestly above GCN, GraphSAGE and GIN (0.69--0.75; ranges overlap, and the differences are indicative only with five seeds). The fine-tuned transformers, run once, follow the same picture: RoBERTa fails comment splitting too (60/90 held-out, external 0.778) and CodeBERT is weaker externally (0.533). LOGWAT's external Benign recall is 0.67 $\pm$ 0.04, so its advantage is relative, not a solved problem. Finally, comment splitting is also solved by a token-level TextCNN with no graph at all, so solving a held-out cell is not by itself evidence that graph structure is needed.

\paragraph{Why LOGWAT fails comment splitting, and what fixes it.}
The failing payloads contain four semantic edges among about 142 sequential and skip edges (2.7\%). Several observations indicate that these edges are diluted by a single shared aggregation rather than absent or uninformative. (i) Deleting LOGWAT's sequential edges, which raises the semantic share to 5.4\%, makes it pass the cell at all five seeds (mean SQLi logit margin $-2.4 \rightarrow +4.4$), so the information is in the model. (ii) HGT, which applies a separate message transformation per edge type, solves the cell at all five seeds (90/90 on the whole matrix). (iii) This is due to the edge types and not to HGT's other differences: two controls with identical architecture, training loop and, for one of them, parameter count (610{,}357), but without relation information (one edge type; three randomly assigned types), solve the cell in 1 of 5 seeds each, versus 5 of 5 (Fisher exact, one-sided, 2/10 versus 5/5: $p=0.007$), and score 80.0 and 81.6 of 90 on the matrix versus 90.0. (iv) Deleting or re-typing only the four semantic edges at inference returns HGT to LOGWAT's failure (margin $+3.96 \rightarrow -2.35$, against $-2.44$ for LOGWAT), whereas the same deletion leaves LOGWAT unchanged. (v) The mechanism is not attention: HGT's first-layer attention on semantic edges is close to uniform, and our earlier ablation that fed the edge type into GATv2's attention logit raised the semantic-edge attention ratio only from 0.76 to 0.82 without changing the decision. Type information must reach the message transformation, not only the attention weights. In this sense LOGWAT's single-relation GATv2 uses only part of the multi-relational structure it constructs, and the earlier statement that semantic edges give the attention mechanism ``a direct logical path'' is not supported by the attention weights, which are no higher on semantic edges than on the other edges (raw ratio 0.76--0.82, not corrected for in-degree).

\paragraph{Bounds on this result.}
The remedy depends on the semantic-edge builder. $E_{sem}$ is created by a keyword-pair matcher, and in our data no benign request carries a semantic edge (0 of 15{,}000 training and 0 of 19{,}293 external benign requests), so ``has a semantic edge'' is a rule of precision 1.0 that a typed channel can read off directly. Typed HGT generalizes to other comment-splitting wordings only where the builder emits an edge (25/30 correct over five seeds, versus 7/30 for LOGWAT); on obfuscations the builder cannot rebuild (for example \texttt{UNI+ON}) it is as blind as LOGWAT (0/15); and it flags benign prose containing keyword pairs such as ``order by \dots\ group by'' (5/30 correct versus 15/30 for LOGWAT). Held-out evidence is also narrower than 90 rows suggest: after masking digits the matrix contains 27 distinct inputs, seven of its nine cells being a single template, and only two cells contain a semantic edge. The additional probes above were written by us for this analysis (six strings per group) and are indicative, not a benchmark. Nor is the matrix a pristine hold-out for any model, because the training-set augmentations were designed with its cells in view (with disjoint field names and a different splitting mechanism). Edge typing is therefore a necessary ingredient of HGT's advantage on this cell, not a general defence against obfuscation.

\paragraph{HGT is unstable on external data.}
HGT's external weighted F1 ranges from 0.33 to 0.82, and three of five seeds label almost every one-token or two-token benign field an attack, mostly XSS. This is not an artifact of the semantic edge (no external benign request has one, and deleting them changes nothing) and not overfitting to the matrix: the controls of equal size do not collapse (0 of 10 versus 3 of 5, $p=0.02$). It arises in a regime the data do not cover: 75\% of external benign requests have one or two tokens, the training set contains no benign or SQLi graph of that size, and all held-out cells have at least ten tokens. For such graphs only the sequential edge type exists, and relabelling those edges changes the prediction in a seed-specific way, which suggests that each seed extrapolates its sequential channel arbitrarily. LOGWAT and the untyped controls are stable there (two-token benign recall 0.94--0.99 at every seed). This does not show that typing itself is unstable, only that HGT with typed channels is fragile outside its training distribution.

\paragraph{Toward combining the two.}
The observations above suggest, as a hypothesis we have not tested, that the two architectures' strengths could be combined: a dedicated aggregation path for $E_{sem}$ (typed messages) alongside GATv2's shared aggregation with self-loops for the dense edges, together with training data that cover short requests, so that the semantic channel is not diluted and the dense channel is not left to extrapolate. Whether this retains LOGWAT's external stability and HGT's held-out behaviour is future work.

\paragraph{Limitations.}
All results use one dataset family in which SQLi and XSS are partly template-generated, five seeds per model with uncorrected significance tests, and checkpoint selection on the test split. The held-out matrix and the additional probes are small and hand-built; the external set is class-imbalanced and dominated by short benign fields. We did not evaluate adaptive adversaries who target the semantic-edge builder, which the results above suggest is the natural point of attack.

%%%%%%%%%%%%%%%%%%%%%%%%%%%%%%%%%%%%%%%%%%%%%%%%%%%%%%%%%%%%%%%%%%%%%%%%%%%%%%%%%%%%%%%%%%%
%% CONCLUSION (replaces \subsection{Conclusion}; keep Future direction, add the hybrid item)
%%%%%%%%%%%%%%%%%%%%%%%%%%%%%%%%%%%%%%%%%%%%%%%%%%%%%%%%%%%%%%%%%%%%%%%%%%%%%%%%%%%%%%%%%%%
\subsection{Conclusion}
We represented HTTP requests as multi-relational token graphs whose semantic edges normalize different obfuscations of the same keyword pair, and classified them with a GATv2 network. Under a like-for-like protocol the standard test split no longer separates architectures, so we compared on a held-out obfuscation matrix and an external dataset. On both, LOGWAT is competitive: on the external dataset it has the best mean weighted F1 (0.796), and on the held-out matrix it matches the single-relation GNNs and a TextCNN, at about 2~ms per request. Its main systematic weakness, comment splitting, is a dilution effect that a relation-typed aggregation removes, and controlled experiments attribute the removal to edge typing; the benefit, however, depends on the semantic-edge builder and comes with false alarms on benign keyword pairs and, for HGT, seed-unstable behaviour on very short requests. Aggregating semantic edges through a dedicated typed path while keeping a stable shared path for dense edges is the most direct next step.

%%%%%%%%%%%%%%%%%%%%%%%%%%%%%%%%%%%%%%%%%%%%%%%%%%%%%%%%%%%%%%%%%%%%%%%%%%%%%%%%%%%%%%%%%%%
%% AUTHOR NOTES -- NOT FOR SUBMISSION
%%%%%%%%%%%%%%%%%%%%%%%%%%%%%%%%%%%%%%%%%%%%%%%%%%%%%%%%%%%%%%%%%%%%%%%%%%%%%%%%%%%%%%%%%%%
% Statements elsewhere in sn-article_v1.tex that contradict the measurements and must change with this text:
%  1. Contribution 4 and Abstract: "97.84% accuracy", "state-of-the-art", "outperforming Bi-LSTM and vanilla GNN by 8-12%": not reproduced (all >= 99.9% on the test split; margins within 0.04 pp).
%  2. Old Table 2 numbers (89.5-97.8%) and latencies (3.7 / 4.1 / 5.8 ...): replaced above. LOGWAT is 1.2-1.6x SLOWER than GCN/GraphSAGE/GIN and 3-4x faster than HGT, not "faster than vanilla GCN/GIN". Independent LOGWAT latency: 2.01 ms CPU / 2.31 ms CUDA (docs/LATENCY.md); 1.79 ms in the same-run comparison.
%  3. Old Discussion bullets ("2.54% over GATv2", "F1 0.98 proves resistance", "faster than GCN/GIN due to sparse adjacency"): remove.
%  4. Section "Visual explainability and case study" (Fig. 6): says GATv2 aggregates UNION and SELECT through the E_sem "shortcut". Measured: LOGWAT's attention on E_sem edges is NOT higher than on other edges (raw ratio 0.76-0.82, not degree-corrected) and it fails the comment-splitting cell 0/50. Rewrite as: the graph builder makes the two obfuscations map to one edge, but a shared aggregation does not exploit it; HGT does.
%  5. Conclusion paragraphs "decisive max-pooling", "effectively neutralized evasion tactics", "robust inductive generalization ... without re-training": not supported at this strength (external Benign recall 0.67).
%  6. "Skip edges provide robustness against junk-code injection": not re-evaluated here; earlier ablation (log section 11) shows case_mixing regresses when skip edges are removed, but verify wording before keeping the claim.
%  7. Dataset description: SQLi pool is 41 templates (16 train-only), 25.2% of the test split duplicates training content; state both in the setup section. Check docs/DATASET.md for the known inaccuracy about "8 payloads per family".
%  8. Do not add HGT to any statement of the form "type-aware attention confirms": the measured mechanism is type-specific messages; attention is near uniform.
